\documentclass[10pt]{article}
\usepackage[letterpaper,margin=0.58in]{geometry}
\usepackage{times}
\usepackage{microtype}
\usepackage{booktabs}
\usepackage{xcolor}
\usepackage{hyperref}
\usepackage{enumitem}
\pagestyle{empty}
\hypersetup{colorlinks=true,linkcolor=black,urlcolor=blue}
\setlength{\parindent}{0pt}
\setlength{\parskip}{3pt}
\setlist[itemize]{leftmargin=1.05em,itemsep=1pt,topsep=1pt}
\definecolor{accent}{RGB}{37,94,142}
\definecolor{soft}{RGB}{245,248,250}
\newcommand{\tagline}[1]{\colorbox{soft}{\parbox{\dimexpr\linewidth-2\fboxsep}{\textcolor{accent}{\textbf{#1}}}}}
\begin{document}

{\Large \textbf{Barrier-Certified Seam Models for Compositional Robot Skills}}\\[-1pt]
{\normalsize One-page collaboration memo for Paper 119}\\[3pt]
\tagline{Core idea: treat skill handoff as a local world/action model: predict whether composition will fail, diagnose why, choose repair/probe/abstain/transition, and update future planning.}

\textbf{Problem.}
Robot skills that work individually can fail when composed. A graph edge may look legal while the terminal state of the first skill falls outside the next skill's attraction basin, crosses a high-energy barrier, or enters a contact mode where the next controller no longer descends. The failure is not only low-level control; it is a representation and prediction problem at the action seam.

\textbf{Method.}
The proposed composer scores candidate handoff $h$ by basin overlap, barrier height, descent continuity, repair cost, and calibrated seam risk:
\[
S(i,j,h)=\alpha O_{\mathrm{basin}}-\beta B_{\mathrm{barrier}}+\gamma D_{\mathrm{descent}}-\lambda C_{\mathrm{repair}}-\kappa R_{\mathrm{seam}}.
\]
It then chooses among \textbf{accept}, \textbf{repair}, \textbf{probe}, \textbf{abstain}, or \textbf{fallback}. The contribution is a seam-level world/action model for behavior composition, not a new low-level controller or universal manipulation policy.

\textbf{Current evidence.}
\begin{center}
\begin{tabular}{lcc}
\toprule
Hard-slice metric & Proposed & Strongest non-oracle predecessor \\
\midrule
Success & 0.80171 & 0.71711 \\
Utility & 0.88827 & 0.65310 \\
Paired utility wins & 10/10 & -- \\
\bottomrule
\end{tabular}
\end{center}
The frozen local suite covers 12 methods, 6 task families, 8 seam regimes, 5 deployment splits, and 10 paired seeds. The proposed method reduces seam failure, barrier violation, damage, risk calibration error, and realized seam breach while improving basin alignment and descent continuity. A diagnostic audit checks exported failure labels, seam decisions, and planner-edge updates over 230,400 local rows.

\textbf{Claim boundary.}
The current claim is bounded: energy-seam certification improves skill composition when basin, barrier, descent, repair, and risk quantities are estimable. The paper should not claim deployment-ready robotics until it has real-robot or accepted high-fidelity validation with shared skill libraries, paired resets, raw logs, videos, and independent baselines.

\textbf{Why this is a natural fit for behavior composition.}
CoStream asks how simple behaviors can compose into complex manipulation. Paper 119 asks when a composed behavior handoff should be trusted, repaired, probed, or avoided. A clean collaboration question is whether a seam model improves robustness for a CoStream-style stack under perturbations, precision assembly, object transfer, and contact-mode transitions.

\textbf{Proposed validation.}
Use the same primitive skills across baselines; log terminal samples, basin estimates, barrier scores, descent diagnostics, predicted risk, accepted/rejected seams, outcomes, and videos. Compare against option graphs, TAMP screens, diffusion stitching, CEM, residual repair, energy heuristics, the previous v4.1 method, and an oracle post hoc upper bound.

\vfill
References for fit: \href{https://haonan16.github.io/}{Haonan Chen homepage}; \href{https://arxiv.org/abs/2606.26423}{CoStream}; \href{https://yilundu.github.io/}{Yilun Du homepage}; \href{https://arxiv.org/abs/2512.05955}{SIMPACT}.

\end{document}
